\section{Introduction}
\label{sec:introduction}

%% DRAFT — rewrite in your own voice before submission.
%%
%% The rubric assesses, in this section and the conclusion, whether *you* understand and
%% can motivate the work. The argument below is sound and every number in it is real, but
%% a grader can tell the difference between an argument you made and one you inherited.
%% Keep the structure, say it the way you would say it.

Albanian is spoken by roughly seven million people and remains under-resourced in natural
language processing. Named entity recognition --- locating mentions of people,
organisations and places in running text --- sits underneath information extraction,
search and question answering, so a language without a dependable NER resource inherits
that gap in everything built on top of it.

Albanian is not unresourced. WikiANN~\cite{bib:wikiann,bib:wikiann-splits} supplies silver
data for the language, and two hand-annotated corpora exist~\cite{bib:albner,bib:kote2025},
the larger of them an order of magnitude bigger than the corpus built here. The gap this
work addresses is therefore not an absence of data but an absence of measurement.

That distinction is what motivates the work, and it can be shown rather than asserted.
WikiANN-sq is built by projecting annotations across Wikipedia links, a method designed for
breadth across 282 languages rather than depth in any one of them, and the Albanian portion
carries the cost of that trade: its median sentence is 6 tokens against 20 in a freshly
sampled Wikipedia corpus, and it contains fragments, leaked wiki markup, and lines
beginning with the redirect directive \texttt{RIDREJTO} presented as sentences
(Section~\ref{sec:results:corpus}). A model scored on such text has been measured on
something other than Albanian prose. This is a limitation of automatic construction, not a
criticism of the method, which achieved what it set out to achieve.

Hand annotation avoids that problem and introduces others, which are usually left
unexamined. Annotators disagree, and how much they disagree bounds how much of a model's
score is signal. Albanian's \emph{X i Y} construction --- a common noun linked to a proper
name, as in \emph{qyteti i Tiranës} --- is genuinely ambiguous about where the name ends,
and a corpus that quietly picks one reading hides a decision that changes its numbers.
Pre-annotating with a model to save effort risks anchoring annotators on the model's
mistakes, converting a labour saving into a correlated error.

This report treats each of those as something to measure. It contributes a hand-verified
Albanian NER corpus on the WikiANN tagset, annotated under an explicit written protocol
with inter-annotator agreement reported rather than assumed; an annotation scheme that
records the full span and the proper-name head of every entity, so the \emph{X i Y}
boundary stays a reporting choice rather than a commitment baked into the data; a
between-subjects experiment in which the same sentences are annotated with and without the
model's suggestions visible, measuring the anchoring effect directly
(Section~\ref{sec:method:annotation}); and an evaluation of published checkpoints and a
fine-tuned model in which every comparison is accompanied by a bootstrap confidence
interval and a paired significance test, so that reported differences are shown to exceed
sampling variation. An active learning comparison is included and, on the calibration data,
returns a null result that the design is reported as being underpowered to resolve
(Section~\ref{sec:results:al}).

Section~\ref{sec:background} covers the necessary background,
Section~\ref{sec:method} documents what was built and how it was measured,
Section~\ref{sec:results} reports the results, and Section~\ref{sec:conclusion} sets out
the limitations.
